\chapter{Limites et perspectives}
\label{chap:limites}

Ce chapitre liste les limites du solveur livré et les pistes d'amélioration documentées dans \texttt{docs/CHOICES.md}.

% =============================================================================
\section{Limites du solveur actuel}
\label{sec:limites:current}
% =============================================================================

\subsection{Régression à très haut nombre de threads}

Le speedup peak observé sur Romeo est 8.23× à 16 threads sur 256×256. Au-delà, régression brutale jusqu'à 0.19× la baseline à 192 threads (5× plus lent que serial).

Trois causes contribuent (analysées en Chapitre~\ref{chap:resultats}) :
\begin{itemize}
    \item Frontières BFS courtes ($\sim 10$ cellules) sur 192 threads.
    \item Fork/join répété ($\sim 50$ µs/niveau × 8000 niveaux).
    \item Contention atomique inter-NUMA (3× plus chère qu'intra-NUMA).
\end{itemize}

L'optim « frontier-threshold » atténue les deux premières (+25\% à 64t, +29\% à 192t), mais la cause (3) reste non-adressée.

\subsection{Plateau smooth\_N3}

Sample \texttt{smooth\_N3} plafonne à 2.23× au peak (4 threads sur 128×128) et régresse dès 8 threads. Cause : peu de tuiles ($L = 12$) + $N = 3$ = peu de parallélisme par cellule (300 ops/cell), frontières BFS courtes (motifs lisses propagent peu).

Pas de fix algorithmique côté intra-attempt sans risquer de régresser binary. Documenté comme fondamental.

\subsection{GPU Kokkos plus lent que CPU sur ce workload}

Le port Kokkos GPU (CUDA via GH200) est fonctionnel mais 8× plus lent que CPU OMP 8t à 128×128. Quatre causes :
\begin{enumerate}
    \item H↔D copies par propagate dominent ($\sim 16$ GB pour 256×256).
    \item Pas assez de parallélisme par niveau BFS (16 896 cores CUDA, niveau de 100 cellules → 99\% oisifs).
    \item Atomics globales du GPU bien plus lentes que CPU local.
    \item Min-entropie reste sur CPU (la phase qui domine ne profite pas du GPU).
\end{enumerate}

% =============================================================================
\section{Pistes d'amélioration}
\label{sec:limites:perspectives}
% =============================================================================

\subsection{NUMA-aware partitioning}

\textbf{Objectif} : adresser la cause (3) de la régression high-thread (contention atomique inter-NUMA).

\textbf{Idée} : partitionner la wave par socket / NUMA node. Chaque thread n'écrit que sur la portion de wave de son NUMA node ; les frontières inter-partitions sont gérées par messages explicites (boîtes aux lettres, pas atomics).

\textbf{Coût d'implémentation} : ~500 lignes, refactor majeur de \texttt{propagate\_tasks}, risque élevé de bugs de race.

\textbf{Bénéfice attendu} : +30-50\% à 64-192 threads (élimination des cache-line ping-pongs cross-NUMA).

\textbf{Statut} : non implémenté. Évalué comme « risque élevé pour bénéfice incertain » dans \texttt{docs/CHOICES.md}.

\subsection{Batch GPU : K solveurs en parallèle dans un seul kernel}

\textbf{Objectif} : adresser les causes (1) et (2) de la performance GPU décevante.

\textbf{Idée} : au lieu de lancer un solveur sur GPU pour une grosse grille, lancer K solveurs indépendants sur K petites grilles, chacun assigné à un SM. Le GPU devient utile quand le data-parallelism massif est exploité.

\textbf{Application} : génération en batch d'une scène avec plusieurs zones procédurales simultanément (un jeu open-world avec 32 régions à générer en parallèle).

\textbf{Statut} : non implémenté. Documenté comme la voie principale pour vraiment exploiter le GPU sur ce workload.

\subsection{Heuristique d'entropie incrémentale}

\textbf{Objectif} : éliminer le re-scan complet de la wave à chaque collapse (qui domine $\sim 85\%$ du temps série).

\textbf{Idée} : maintenir une priority queue des cellules par entropie, mise à jour incrémentale à chaque modification. Lazy invalidation : quand une cellule est touchée par propagate, on la marque \textit{stale} dans la PQ, on la re-priorise au prochain pop.

\textbf{Bénéfice attendu} : la sélection passe de $O(r \cdot c \cdot L)$ par collapse à $O(\log(r \cdot c))$ amorti. Pour 256×256 ($\sim 65k$ cellules), facteur $\sim 1000$ sur le coût de la sélection.

\textbf{Coût d'implémentation} : modéré. Une PQ thread-safe est délicate ; une variante plus simple est de maintenir un \textit{tournament tree} indexé par cell\_id.

\textbf{Statut} : non implémenté. C'est probablement la piste avec le meilleur ratio bénéfice/risque.

\subsection{SIMD pour les opérations bitset}

\textbf{Objectif} : accélérer la version série pour les samples avec $L > 64$ (bitsets multi-mots).

\textbf{Idée} : remplacer les boucles \texttt{for (i ; i < n ; ++i) words[i] op= other[i]} par des intrinsics AVX-512, traitant 8 mots de 64 bits par cycle.

\textbf{Bénéfice attendu} : facteur 4-8 sur les AND/OR multi-mots. Effet limité car la version série est déjà O(words\_per\_cell) avec le compilateur en \texttt{-O3 -march=native} qui auto-vectorise dans certains cas.

\textbf{Coût d'implémentation} : faible (juste un AVX-512 wrapper). Risque de portabilité (pas tous les CPU ont AVX-512, ARM Neoverse-V2 a SVE).

\textbf{Statut} : non implémenté. Bénéfice incertain.

\subsection{Symétries dans les règles d'adjacence}

\textbf{Objectif} : extension naturelle de l'option \texttt{--symmetries} actuelle. Au lieu d'étendre seulement le tile set, exploiter les symétries dans la \textit{construction des règles} pour réduire la mémoire.

\textbf{Idée} : si $t' = R_{90}(t)$, alors $\text{allowed}(t', dx, dy) = \text{rotate90}(\text{allowed}(t, dy, -dx))$. On peut stocker uniquement les règles pour les tuiles \textit{canoniques} (un représentant par classe d'équivalence) et déduire les autres à la demande.

\textbf{Bénéfice} : mémoire des règles divisée par jusqu'à 8.

\textbf{Statut} : non implémenté. Niche.

% =============================================================================
\section{Limites de l'évaluation}
\label{sec:limites:evaluation}
% =============================================================================

\subsection{Une seule architecture HPC testée}

Tous les benchmarks « HPC » sont sur Romeo (EPYC 9654). Nous n'avons pas testé sur :
\begin{itemize}
    \item Intel Xeon (architectures différentes, latences mémoire différentes).
    \item AMD Genoa-X (V-Cache, pourrait changer le profil cache).
    \item ARM Graviton (cloud AWS).
\end{itemize}

Les conclusions sur le scaling (peak à 16 threads sur 256×256, régression au-delà) sont spécifiques à EPYC 9654 et probablement transférables à d'autres EPYC, mais à valider.

\subsection{Tailles limitées}

Le sweep principal monte à 256×256. Une exécution à 512×512 a été tentée mais timeout (4h). Les conclusions sur le scaling ne peuvent pas être extrapolées au-delà sans mesure.

\subsection{Variance non quantifiée}

Nos mesures sont des médianes de 3 ou 11 runs. Nous ne reportons pas d'intervalle de confiance ou d'écart-type. Pour les A/B critiques (frontier-threshold), une analyse statistique formelle (test de Mann-Whitney) serait plus rigoureuse.

